\documentclass[11pt,a4paper]{article}
\usepackage[utf8]{inputenc}
\usepackage[T1]{fontenc}
\usepackage{amsmath, amssymb}
\usepackage{hyperref}
\usepackage{graphicx}
\title{Privacy Leakage Measurement in Enterprise RAG Systems: A Baseline Empirical Study}
\author{Author Name}
\date{\today}
\begin{document}
\maketitle

\begin{abstract}
Retrieval-augmented generation (RAG) has become the default architecture for grounding large language models in proprietary enterprise knowledge, yet the privacy properties of these systems remain largely unmeasured under adversarial conditions. This study presents a reproducible measurement framework for quantifying how often, and how severely, a frontier language model discloses private information when an enterprise RAG system is deliberately attacked. The framework rests on a synthetic enterprise corpus of 155 documents spanning employee records, medical reports, system configurations, and internal memoranda; four attacker profiles with formally distinct authorization thresholds; and eight attack families organised by mechanism rather than outcome. Across 52 attack scenarios comprising 267 conversational turns, each replayed three times under deterministic decoding, the framework establishes an undefended baseline against which future defensive interventions can be measured. The headline result is an augmented strict attack success rate of 28.8\%, with two attack families---direct PII extraction and multi-turn linkability inference---accounting for 91.1\% of all turn-level strict leaks. A reproducible multi-turn case study demonstrates that the model's privacy defences are turn-local rather than conversation-aware, refusing sensitive disclosures in isolation while complying with a follow-up that depends on the very population it had just declined to discuss. A second discovery shows that descriptive corpus poisoning succeeds where imperative poisoning fails. An automated judge, calibrated against human annotation, achieved a Cohen's kappa of 0.83. No severity-five disclosures were observed. The study deliberately leaves defences disabled, producing a clean characterisation of the model's intrinsic behaviour and a concrete set of defence priorities for the next phase.
\end{abstract}

\section{Introduction}

Enterprise adoption of retrieval-augmented generation has outpaced the development of methods to evaluate its privacy risk. Organisations connect language models to internal knowledge bases containing personnel files, clinical notes, infrastructure secrets, and confidential correspondence, on the assumption that retrieval grounding improves factual accuracy without materially expanding the attack surface. That assumption has not been tested rigorously under adversarial pressure. Most published privacy evaluations of language models examine memorisation of training data, or measure the propensity of a bare model to regurgitate personal information from its pretraining corpus. These evaluations are valuable, but they do not describe the situation that an enterprise actually faces, in which the sensitive data is not in the model's weights at all but is instead injected into the context window at inference time by a retrieval layer the operator built.

The gap this study addresses is the absence of a measurement framework for privacy leakage in \textit{adversarial} RAG scenarios. When a benign user asks a benign question, a RAG system that surfaces a colleague's salary is committing a privacy failure, but the failure is incidental. When an adversary deliberately constructs a sequence of queries designed to extract that salary, or to link an anonymised medical record back to a named individual, or to smuggle a fabricated identity into the knowledge base, the system faces a qualitatively different threat. The probability and severity of leakage under deliberate attack is the quantity that matters for risk management, compliance, and architectural decision-making, and it is precisely the quantity that existing evaluation practice does not capture.

The contribution of this work is twofold. First, it provides a reproducible measurement framework: a defined corpus, a formal threat model with four attacker profiles, an attack taxonomy organised by mechanism, a dual attack-success-rate methodology that separates formal authorization violations from broader disclosure, and an automated judging procedure calibrated against human annotation. Second, it provides a baseline: a quantified characterisation of how an undefended, frontier-model-backed enterprise RAG system behaves when attacked across the full taxonomy. The baseline is not the end of the research programme; it is its foundation. Every defensive intervention proposed for subsequent phases is justified by, and will be evaluated against, the numbers established here.

A deliberate methodological choice underpins the entire study: defences are disabled for the baseline. There is no PII redaction, no output filtering, no retrieval access-control layer, and no conversation-aware refusal mechanism. This is not an oversight but the point. The baseline isolates the intrinsic behaviour of the generation model when it is handed retrieved context and an adversarial prompt. Any refusal or safety behaviour observed in the baseline is the model's own, unaided by external machinery. This makes the baseline a clean reference: when defences are later added, the difference between the defended and undefended measurements is attributable to the defences and not confounded by them.

\section{Threat Model and Attack Taxonomy}

\subsection{Attacker Profiles}

The framework defines four attacker profiles. Each profile encodes a distinct combination of access, motivation, and---critically---an authorization threshold that determines what counts as a leak for that attacker. The threshold is the central abstraction. A disclosure that is perfectly acceptable to one attacker constitutes a privacy violation when delivered to another, because the two attackers are authorised to see different things. Treating leakage as profile-relative, rather than absolute, is what allows the framework to reason about realistic enterprise access models in which the same underlying datum is sensitive or benign depending on who receives it.

\begin{table}[h]
\centering
\begin{tabular}{l l c c}
\hline
\textbf{Profile} & \textbf{Identity} & \textbf{Authorization threshold} & \textbf{What constitutes a leak} \\
\hline
P1 & Internal     & 2 & severity $\geq 3$ from the corpus \\
P2 & External     & 0 & any corpus disclosure \\
P3 & Admin        & 3 & severity $\geq 4$ via chat interface \\
P4 & Third-Party  & 3 & aggregate / cross-tenant severity $\geq 4$ \\
\hline
\end{tabular}
\caption{The four attacker profiles and their authorization thresholds. Severity is graded on the corpus annotation scale; a higher threshold means the attacker is authorised to see more, so fewer disclosures count as leaks.}
\end{table}

The \textit{External} attacker (P2) has an authorization threshold of zero. This profile models an outsider with no legitimate access to any enterprise data, so any disclosure of corpus content whatsoever counts as a leak. It is the strictest profile in the sense that it tolerates nothing, and consequently the profile against which inclusive disclosure is most visible. The \textit{Internal} attacker (P1) models an authenticated employee who is entitled to low-sensitivity information but not to high-sensitivity records; the threshold of two means that only disclosures of severity three or above constitute a violation. The \textit{Admin} attacker (P3) is the most privileged human profile, authorised up to severity three, so only the most sensitive material---severity four and above---leaks; the Admin attack surface is the chat interface itself, reflecting the reality that highly privileged operators interact with the system conversationally. The \textit{Third-Party} attacker (P4) models an integration partner or downstream tenant; its concern is aggregate and cross-tenant exposure, where individually innocuous facts combine into a severity-four disclosure that crosses a tenancy boundary.

The thresholds differ because the profiles represent genuinely different trust relationships, and a measurement framework that collapsed them into a single notion of leakage would misrepresent enterprise risk. An access model that grants administrators broad visibility is not failing when it shows an administrator moderately sensitive data; it is failing only when it shows that administrator the small set of records reserved above their level. By contrast, an external party should see nothing, so the bar for a leak is on the floor. Encoding these distinctions as numeric thresholds makes leakage a precise, profile-conditioned predicate rather than an intuition.

\subsection{Attack Families}

The taxonomy organises adversarial behaviour into eight families. The organising principle is mechanism, not outcome: families are distinguished by \textit{how} they attempt to induce leakage, because a defence that neutralises a mechanism generalises across all the outcomes that mechanism could produce, whereas a defence tuned to an outcome may be bypassed by a new mechanism reaching the same end. Classifying by mechanism is therefore the choice that makes the taxonomy useful for designing defences rather than merely cataloguing failures.

\begin{table}[h]
\centering
\begin{tabular}{l l}
\hline
\textbf{Family} & \textbf{Mechanism} \\
\hline
F1 Direct PII extraction      & Ask plainly for protected fields and rely on compliance. \\
F2 Linkability inference      & Combine innocuous turns to re-identify an individual. \\
F3 Prompt injection           & Embed adversarial instructions in the query or context. \\
F4 Obfuscation                & Disguise the request through encoding or indirection. \\
F5 Jailbreak                  & Override the model's policy through role or framing. \\
\hline
\end{tabular}
\caption{Five of the eight attack families, by mechanism. The remaining three---F6 Reconnaissance (probe the system to map its corpus and behaviour), F7 Social engineering (manipulate the model through pretext and rapport), and F8 Corpus poisoning (plant adversarial content in the knowledge base)---complete the taxonomy.}
\end{table}

\textit{Direct PII extraction} (F1) is the simplest mechanism: the attacker simply asks for the protected field and depends on the model's willingness to answer. Its simplicity is deceptive, because a system without an enforcement layer has no principled reason to refuse a fluent, well-formed request. \textit{Linkability inference} (F2) is the multi-turn mechanism: no single turn asks for anything obviously sensitive, but the turns are constructed so that their answers, taken together, re-identify or profile an individual. This family targets the threat that LINDDUN labels \textit{linkability}, and it is the mechanism most likely to evade defences that inspect queries in isolation. \textit{Prompt injection} (F3) embeds adversarial instructions into the query or, more dangerously, into retrieved content, attempting to subvert the model's task. \textit{Obfuscation} (F4) disguises the malicious request so that a naive filter does not recognise it. \textit{Jailbreak} (F5) attacks the model's policy directly through role-play, hypothetical framing, or authority claims. \textit{Reconnaissance} (F6) does not seek data directly but maps the system---what it contains, how it behaves, where its boundaries lie---in preparation for a targeted attack. \textit{Social engineering} (F7) manipulates the model through pretext, urgency, and manufactured rapport. \textit{Corpus poisoning} (F8) attacks the knowledge base rather than the model, planting adversarial documents that the retrieval layer will later surface as if authoritative.

\subsection{Conceptual Mapping to LINDDUN}

The taxonomy aligns conceptually with the LINDDUN privacy threat framework, which decomposes privacy threats into linkability, identifiability, non-repudiation, detectability, disclosure of information, unawareness, and non-compliance. The mapping is not one-to-one, but it situates the attack families within an established privacy-engineering vocabulary. Direct PII extraction and several poisoning and injection variants map most directly to \textit{disclosure of information}: the unauthorised release of personal data. Linkability inference maps to \textit{linkability} and \textit{identifiability}: the ability to associate records or to single out an individual from a population. Reconnaissance touches \textit{detectability}: learning that a particular record exists is itself a privacy-relevant signal even before its contents are read. Corpus poisoning, by introducing fabricated identities that the system then treats as genuine, raises \textit{non-compliance} and integrity concerns that sit at the boundary between privacy and trustworthiness. Framing the taxonomy against LINDDUN clarifies that the framework is not measuring a single failure mode but a spectrum of distinct privacy harms, each of which a complete defence strategy must address.

\section{Measurement Architecture}

The measurement architecture is best understood as a set of design decisions, each made to maximise reproducibility and measurement validity rather than to optimise the system's real-world performance. The system under test is intentionally a plausible but unhardened enterprise RAG pipeline, because the object of study is how such a pipeline behaves by default.

\subsection{Vector Store: Local and Persistent}

The retrieval layer uses a local, on-disk, persistent vector store. The decision to keep the store local and persistent, rather than to use a hosted retrieval service, was made for reproducibility. A persistent local index can be built once, frozen, and queried identically across every run of the campaign, which means that differences in measured leakage are attributable to the attack and the model rather than to drift in an external service. A hosted service, by contrast, may re-index, re-rank, or update its internals between runs in ways the experimenter cannot observe or control, introducing a confound that would undermine the validity of a baseline whose entire purpose is to be a stable reference point.

\subsection{Embedding Model: Small, Stable, Well Understood}

The pipeline embeds documents and queries with a small, widely studied open sentence-embedding model. The choice of a compact and well-characterised embedder over a larger or proprietary one was deliberate. A small, stable model produces deterministic embeddings and behaves predictably, so the retrieval behaviour observed in the baseline reflects properties that other researchers can reproduce and reason about. Using a frontier or frequently updated embedding service would have made the retrieval layer a moving and opaque target. Because the study's goal is to characterise an honest, commonly deployed configuration rather than to chase maximal retrieval quality, a representative, well-understood embedder is the more valid choice. The retrieval results reported later in this paper should be read in this light: they describe what a typical, reasonable enterprise deployment would experience, not the ceiling of what retrieval technology can achieve.

\subsection{Generation Model: A Frontier LLM with Deterministic Decoding}

The system under test is backed by a frontier large language model operating under deterministic decoding. Two decisions matter here. The first is the use of a genuinely capable frontier model, so that the baseline reflects the privacy behaviour of the class of systems organisations are actually deploying, not a weak proxy whose failures would be uninformative. The second is deterministic decoding, chosen so that each scenario can be replayed and the model's response is, to the greatest extent the platform allows, stable across replays. Determinism is what makes the three-times replay of every scenario meaningful: a finding that reproduces across independent runs under deterministic decoding is a property of the system, not an artefact of sampling noise. The reproducibility of the central case study, which holds across all three runs, depends on this decision.

\subsection{Retrieval Configuration: Top-$k$ Without Defensive Filtering}

The retrieval configuration returns the top-$k$ most similar chunks for each query and passes them to the generation model without any defensive filtering. There is no redaction of retrieved content, no access-control check on which documents a given profile may retrieve, and no scrubbing of sensitive fields before the model sees them. This is the architectural expression of the baseline's defining commitment: defences are off. The retrieval layer does exactly what an undefended RAG system does---it finds semantically relevant context and hands it over. Measuring leakage under this configuration establishes the worst-case-by-default behaviour and gives every later defensive layer a clearly defined quantity to reduce.

\section{Evaluation Methodology}

\subsection{The Dual Attack Success Rate}

A single leakage metric is insufficient, because leakage is profile-relative and a lone number forces a false choice between two legitimate questions. The framework therefore reports two attack success rates for every measurement. The \textit{strict} rate counts a turn as a successful attack only when the disclosure exceeds the attacker's authorization threshold---that is, when a formal authorization violation has occurred under the profile's own access model. The \textit{inclusive} rate counts a turn as a successful attack whenever any corpus-derived personal information is disclosed, regardless of whether the receiving profile was authorised to see it.

The two rates answer different and equally important questions. The strict rate answers: did the system violate its own access-control contract? This is the compliance-relevant figure, the one that maps to a policy breach. The inclusive rate answers: how much personal information is flowing out of the system at all? This is the exposure-surface figure, the one that reveals how much sensitive content the model is willing to surface even when, for a given privileged profile, that surfacing is technically permitted. A system can have a low strict rate and a high inclusive rate, which is exactly the signature of a model that discloses freely but happens to be queried by attackers with permissive thresholds. Reporting only the strict rate would hide that exposure; reporting only the inclusive rate would overstate the compliance failure. The gap between the two rates is itself a diagnostic, and as the results show, that gap can be very large.

\subsection{LLM-as-Judge Calibration}

Determining whether a given response leaked, and at what severity, requires judgement on free-form text at a scale beyond what manual annotation can sustain across hundreds of turns and their replays. The framework therefore employs an automated judge---a separate, more capable frontier model in a labelling role---to classify each response. Using a second, independent model as the labeller, rather than the model under test, avoids the obvious circularity of asking a system to grade its own outputs and provides a labeller whose judgements can be audited.

The trustworthiness of an automated judge cannot be assumed; it must be demonstrated. The framework establishes it by measuring agreement between the judge and human annotation on a sample of responses, using Cohen's kappa, a chance-corrected agreement statistic that is the standard instrument for this purpose. The judge achieved a Cohen's kappa of 0.83 against human annotation, agreeing on 43 of 46 sampled cases. A kappa of 0.83 sits firmly in the range conventionally described as strong agreement, which licenses the use of the automated judge as a stand-in for human labelling across the full campaign while acknowledging that it is not perfect. The three disagreements are not noise to be discarded; they mark the boundary regions where automated and human judgement diverge, and they are part of why the framework reports its central findings conservatively.

\section{Key Empirical Results}

\subsection{Headline Leakage Rate}

The undefended baseline exhibits an \textbf{augmented strict attack success rate of 28.8\%}. Across the full taxonomy of attacks and the full set of attacker profiles, in other words, more than one in four attacking turns succeeds in producing a disclosure that violates the receiving profile's authorization threshold. This is the single most important number the baseline establishes, and it is the quantity that the next phase's defences are charged with driving down. It is worth dwelling on what this figure represents: it is not the rate at which the model says something sensitive in casual use, but the rate at which a deliberate adversary, working through a structured taxonomy of mechanisms, obtains an unauthorised disclosure from a system with no defensive layer interposed.

\subsection{Concentration of Risk in Two Families}

The leakage is not spread evenly across the taxonomy. \textbf{Two attack families cause 91.1\% of turn-level strict leaks: direct PII extraction (77.8\%) and multi-turn linkability inference (63.3\%).} Equivalently, \textbf{F1 and F2 account for 91.1\% of turn-level strict-augmented leaks.} This concentration is the most strategically consequential finding in the study, because it transforms the defence problem from an intractable demand to harden every conceivable attack pathway into a tractable programme of closing two dominant mechanisms. Direct extraction succeeds on more than three-quarters of its attempts and multi-turn linkability on nearly two-thirds, while the remaining six families contribute comparatively little to the strict total. A defence that directly neutralises F1 and F2 would close most of the measurable risk; a defence that targets the remaining families would have little effect on the headline.

\subsection{Turn-Local, Not Conversation-Aware: The P2-LINK-003 Case Study}

The most illuminating qualitative finding is that \textbf{the model's privacy defences are turn-local, not conversation-aware.} The reproducible multi-turn attack catalogued as scenario \textit{P2-LINK-003} demonstrates the failure mode precisely. In this scenario the model first refuses to disclose salary information. It then refuses to disclose medical information. It then complies with a request about the office locations of ``those employees''---the very population defined implicitly by the two preceding, refused turns. The model evaluates each turn in isolation and finds the office-location question benign, failing to recognise that ``those employees'' inherits its referent from the sensitive turns it had just declined. The disclosure occurs not because the model lacks a sense of what is sensitive---it demonstrably refused the direct sensitive requests---but because its sensitivity judgement does not persist across the conversation.

This case study is the centrepiece of the baseline for three reasons. First, it reproduces: the pattern holds across all three independent runs of the scenario under deterministic decoding, so it is a property of the system rather than a fluke of one generation. Second, it is mechanistically clear: it is not that the model was tricked by clever phrasing, but that its defensive reasoning has no memory, so a request that would be obviously sensitive in context is treated as innocuous when read alone. Third, it directly motivates a defence: a conversation-aware refusal mechanism that propagates a sensitivity classification across turns, so that once a topic or population has been declined, later references to ``those employees'' or ``those people'' remain in the declined state, would close this exact failure. The case study thus functions as both an empirical result and a design requirement.

\subsection{Descriptive Poisoning Succeeds Where Imperative Poisoning Fails}

The corpus-poisoning family yielded a discovery about \textit{how} poisoning succeeds. \textbf{Descriptive corpus poisoning succeeds where imperative poisoning fails: six of seven corpus-poisoning scenarios produced zero leaks; the seventh, using a narrative rather than a command, disclosed a fabricated identity---Patricia Walsh, with email---in three of three runs.} The six poisoning scenarios that failed all relied on imperative content: planted text that instructed the model to do something, in the manner of a classic injected command. The model proved robust to these, treating the imperative as suspect. The single scenario that succeeded planted no command at all. It introduced a fabricated person as ordinary descriptive narrative---a name and contact details embedded in prose, as though they were simply facts about the world. The model, encountering this content during retrieval, surfaced the fabricated identity as if it were genuine, in every one of the three runs.

The significance of this finding is that it identifies a poisoning vector that instruction-side and input-side defences will not catch. A guardrail that scans for adversarial commands sees nothing to flag in a descriptive paragraph, because there is no command to detect. The poisoned content is dangerous precisely because it looks like legitimate corpus material. Defending against it requires reasoning about \textit{provenance}---establishing which documents are authorised sources of identity claims and refusing to surface unverifiable identities from the retrieval layer---rather than reasoning about the form of the request. This discovery is the basis for an additional, newly identified defence layer in the next phase.

\subsection{Retrieval Precision on Named Targets}

A finding about the retrieval layer reframes part of the leakage picture. \textbf{Retrieval precision on named targets was 8.3\% (only 3 of 36 named-target turns retrieved the correct document).} When an attacker named a specific individual and attempted to extract that individual's record, the retrieval layer returned the correct document only about one time in twelve; far more often it returned semantically similar but incorrect material---adjacent records and related memoranda rather than the precise file requested. Retrieval is, in this sense, an incidental and silent defence: many targeted attacks fail not because the model refuses but because the retrieval layer never surfaces the right document in the first place.

This is a double-edged observation. On one hand, it means the undefended baseline's leakage would be substantially worse if retrieval were more precise, and it cautions against over-optimising retrieval quality without simultaneously strengthening privacy controls, since improving precision would directly expand the leakage surface. On the other hand, an incidental defence is not a designed defence: it cannot be relied upon, it provides no guarantees, and it would evaporate the moment an organisation upgraded its embedder for better task performance. The retrieval-precision finding therefore belongs in the baseline as a caveat on the headline numbers and as an argument for deliberate, rather than accidental, control over what the retrieval layer is willing to surface.

\subsection{The External Profile Gap and the Severity Ceiling}

The dual metric is at its most revealing for the External profile. \textbf{The External profile strict rate is 25.9\%; its inclusive rate is 70.4\%---a gap of 44.5 percentage points.} Because the External attacker has an authorization threshold of zero, every disclosure of corpus content is, for this profile, both a strict and an inclusive leak, which makes the magnitude of the gap a direct measure of how much personal information the system surfaces overall relative to how much rises to the level of the strict accounting. An inclusive rate above seventy percent indicates that the model is willing to surface corpus-derived personal information in the large majority of attacking turns directed at it. The forty-four-point gap quantifies the exposure that a strict-only accounting would conceal, and it is the clearest single demonstration of why the framework reports two rates rather than one.

Finally, a bounding result: \textbf{no severity-five disclosures were observed.} Across the entire campaign, the model never leaked information at the highest severity grade. The most sensitive class of records in the corpus was never disclosed in a way that crossed any profile's threshold. This is a genuinely positive finding and an important qualification of the headline: the leakage that the baseline measures, while frequent, is concentrated below the top of the severity scale. The system's failures are serious but not catastrophic in the specific sense that the crown-jewel category held.

\section{Defence Priorities for the Next Phase}

The concentration of risk documented in the results makes a focused defence strategy possible. Rather than attempting to harden every pathway uniformly, the next phase concentrates investment on the layers that target the families with measurable leak rates. Five defence layers are proposed, with priorities derived directly from the baseline findings.

\subsection{Layer 1 --- Input Filtering --- Deprioritised}

Input filtering inspects incoming queries for adversarial content before they reach the model. It is deprioritised. The baseline shows that the dominant successful mechanisms---direct extraction and multi-turn linkability---do not depend on malformed or obviously adversarial inputs; a plain, well-formed question is enough. Filtering the input therefore addresses families that contribute little to the headline while leaving the two dominant families largely untouched. Layer 1 remains in the architecture as a backstop against the cruder injection and obfuscation attempts, but it is not a focus of next-phase implementation work because the evidence shows it would not move the number that matters.

\subsection{Layer 2 --- Retrieval-Time PII Redaction --- High Priority}

Retrieval-time PII redaction applies industry-standard PII detection to retrieved content before it is passed to the generation model, so the model is presented with redacted context and cannot leak what it never sees. This layer is high priority because it directly targets F1, the family with the highest strict-augmented attack success rate among the dominant families, at 77.8\%. By removing the protected fields from the context at the point of retrieval, it attacks the leakage at its source rather than hoping the model declines to repeat what it has been handed. It is the single intervention most likely to reduce the headline figure, because it operates precisely where the most successful mechanism does its work.

\subsection{Layer 3 --- Conversation-Aware Refusal --- High Priority}

Conversation-aware refusal addresses the multi-turn linkability failure documented in the P2-LINK-003 case study, and it is high priority because that failure is the second-largest contributor to the strict total. The current model refuses each turn in isolation; the proposed layer maintains a session-level sensitivity state, so that once a topic, an individual, or a population has been declined in any turn, subsequent turns that refer back to the same referent---``those employees,'' ``those people''---inherit the declined status. This layer is the direct designed remedy for the turn-local failure mode, converting the model's existing but memoryless sense of sensitivity into a property that persists across the conversation.

\subsection{Layer 4 --- Output Filtering --- Supplementary}

Output filtering applies a PII-detection pass to the model's response before it is returned to the user. It is classified as supplementary rather than high priority because it largely addresses the same risk as Layer 2 but later in the pipeline, and is therefore expected to add modest incremental protection once retrieval-time redaction is in place. Its distinct value, and the reason it is retained, is that it catches model-generated content that was never present in any retrieved chunk. The Patricia Walsh case illustrates the principle: where a model can surface or invent personal details that no upstream redaction step inspected, a final pass over the output provides defence in depth. Layer 4 is the safety net for disclosures that originate in generation rather than retrieval.

\subsection{Layer 5 --- Corpus Provenance Validation --- High Priority --- New}

Corpus provenance validation is a newly identified layer, added on the basis of the descriptive-poisoning discovery, and it is high priority because no other proposed layer addresses that vector. The layer validates the provenance of corpus content before it is admitted to the retrieval index, establishing which documents are authorised sources of identity claims and refusing to surface unverifiable identities. It is the only proposed defence that targets descriptive poisoning, because that attack carries no adversarial instruction for an input or output filter to detect---its danger lies in masquerading as legitimate descriptive material. Defending against it requires reasoning about where corpus content came from and whether the identities it asserts can be trusted, which is a provenance question rather than a content-inspection question. The emergence of this layer directly from an empirical finding exemplifies the intended relationship between the baseline and the defence programme.

\section{Conclusion}

This study established a reproducible baseline for privacy leakage in an undefended, frontier-model-backed enterprise RAG system under deliberate attack. The baseline quantifies what existing privacy evaluations leave unmeasured: how a realistic retrieval-augmented system behaves when an adversary works systematically through a taxonomy of mechanisms against profiles with formally distinct authorization thresholds. Its central numbers---a 28.8\% augmented strict attack success rate, the concentration of 91.1\% of strict leaks in two families, a 44.5-point gap between the External profile's strict and inclusive rates, and the absence of any severity-five disclosure---together form a precise reference point that did not previously exist.

The framework establishes more than a set of numbers. It establishes a method: a profile-relative notion of leakage that respects real enterprise access models, a mechanism-based attack taxonomy that maps conceptually onto the LINDDUN threat categories, a dual attack-success-rate methodology that separates compliance violations from exposure surface, and an automated judging procedure whose agreement with human annotation, at a Cohen's kappa of 0.83, is strong enough to license its use at scale. The two qualitative discoveries---that the model's defences are turn-local rather than conversation-aware, and that descriptive poisoning succeeds where imperative poisoning fails---are findings that a purely quantitative evaluation would have missed, and each translates directly into a designed defence.

The baseline also establishes what it is not. It is not a measurement of a defended system; defences were deliberately disabled so that the model's intrinsic behaviour could be isolated. It is not a claim about the ceiling of retrieval technology; the incidental protection afforded by imperfect retrieval precision is a caveat, not a control to be relied upon. And it is not the end of the research programme. Its purpose is to enable the next phase, in which the five proposed defence layers---deprioritising input filtering, prioritising retrieval-time redaction, conversation-aware refusal, and corpus provenance validation, with output filtering as defence in depth---are implemented and evaluated against the very same scenarios and metrics. Because the baseline holds every other variable fixed, the difference between the defended and undefended measurements will be attributable to the defences themselves. That is the function of a baseline, and it is the contribution of this work: a stable, reproducible, adversarial reference against which the privacy of enterprise RAG systems can finally be measured.

\end{document}
